% Introduction and framing draft -- Project Sidewalk disagreement study.
% Pairs with results_draft.tex. Numbers there are verified against the data;
% the citation keys below still need entries in the .bib.

\title{Two Kinds of Disagreement: What Quality Control Removes from
Crowdsourced Urban Accessibility Data}

\begin{abstract}
Cities increasingly rely on crowdsourced assessments of the built environment to
decide where to spend accessibility budgets. Those assessments are cleaned
before use, and cleaning treats reviewer disagreement as noise. We show that it
is not one thing. Across 811{,}368 labels from twelve Project Sidewalk
deployments, reviewers fail to agree in two distinct ways: they split into
opposing camps about whether a condition meets the threshold, or they
individually decline to judge because they cannot tell from the image. The two
are only weakly related, and they draw on different sources---disagreement is a
property of the sidewalk by a factor of five over the reviewer, whereas
uncertainty is barely so. Neither resolves with additional review. Yet the
platform's published quality criterion retains 54.5\% of reviewed labels overall
while retaining only 8.6\% of split labels and 6.8\% of uncertain ones,
discarding over 90\% of both in all twelve cities. The record of where human
judgement about accessibility diverges is being removed before anyone can read
it. We argue that contested labels should be routed rather than deleted, and
that the two kinds require different destinations: better imagery for one, a
policy decision for the other.
\end{abstract}

\section{Introduction}

Whether a stretch of pavement is passable in a wheelchair is not always a
question with an answer. A cracked slab is a hazard to one person and an
inconvenience to another; whether a corner \emph{needs} a curb ramp depends on
what counts as a pedestrian route. These are judgement calls, and cities now
make them at scale by asking volunteers.

Project Sidewalk~\parencite{saha2019projectsidewalk} is among the largest such
efforts: volunteers examine street-level imagery and mark accessibility
problems---missing curb ramps, surface defects, obstructions---and other
volunteers then review each mark as \emph{agree}, \emph{disagree}, or
\emph{not sure}. The resulting labels feed accessibility audits, route planning
for wheelchair users, and municipal remediation priorities. Before any of that,
the data is cleaned: labels that failed to attract sufficient agreement are
filtered out~\parencite{li2024labelaid}.

That filtering step embeds an assumption worth examining. It treats a label that
reviewers argued about as defective---as though disagreement indicates a
mistake by the original labeller, to be discarded like a mis-click. But
disagreement among careful reviewers can also indicate that the underlying
question is genuinely contested, which is a different thing entirely and carries
different consequences. If reviewers split over whether a particular crack rises
to the level of a ``surface problem,'' that split is evidence about where a
city's threshold actually sits, and about how much its accessibility standards
are shared. Deleting it does not resolve the ambiguity; it conceals it.

We further argue that disagreement is not one phenomenon. A reviewer can fail to
endorse a label in two ways that look similar in aggregate and mean opposite
things. Reviewers can \emph{split}: some confidently agree, others confidently
disagree, and the disagreement is about where the standard lies. Or reviewers can
be \emph{unsure}: nobody takes a position, because the image does not settle the
question. The first is a disagreement about values, and more reviewing will not
converge it. The second is a disagreement about evidence, which better imagery
could in principle fix. A pipeline that cannot tell them apart cannot route
either one correctly.

Project Sidewalk is unusual in making this separable at all. Most crowdsourcing
interfaces offer only agreement or disagreement; the explicit ``not sure''
option, recorded per reviewer alongside reviewer identity and position in the
review sequence, is what makes the decomposition possible.

Our contributions are three:

\begin{itemize}
\setlength\itemsep{0.1em}
\item An \textbf{empirical decomposition of reviewer disagreement} into a
  standards component and an evidence component, applied to 746{,}975
  individual review decisions across twelve city deployments. The two are only
  weakly related ($\rho = 0.14$) and differ sharply in source: splitting is a
  property of the label rather than the reviewer by roughly five to one, whereas
  for uncertainty the ratio is close to parity (Section~\ref{sec:results}).
\item Evidence that \textbf{neither kind resolves with additional review}.
  Comparing a label's first three reviews against its later ones, within the
  same label and adjusted for who reviewed it, neither form of disagreement
  declines. Whatever is producing these splits, more eyes do not settle it.
\item A \textbf{measurement of what quality control removes}. The published
  filter retains over half of all reviewed labels but under a tenth of the
  contested ones, in every one of the twelve cities. We give the resulting
  design implication: contested labels are a signal a city needs, and should be
  routed by kind rather than filtered by consensus.
\end{itemize}

A methodological point runs alongside these. Each of the findings above changed
sign or magnitude when we controlled properly for who reviewed a label and for
which city it came from. Reviewer ``not sure'' rates range from 0\% to 70\%,
and deployments differ enough that estimates pooled across them absorb
between-city variation into reviewer effects. We report the uncontrolled
analyses beside the controlled ones, because in each case the uncontrolled
version told a tidier story than the data supports, and because anyone repeating
this work on similar data will meet the same trap.

% ---------------------------------------------------------------------------
% Framing notes (not for the submitted text)
%
% What this paper is NOT claiming:
%   - not that Project Sidewalk's filter is badly designed. It does what it was
%     built to do, which is remove low-quality labels. The claim is about what
%     else goes with them, and that no separate channel preserves it.
%   - not that contested labels are correct. They are contested; correctness is
%     the wrong frame, which is the point.
%   - not that "unsure" is purely an image-quality signal. Within cities it is
%     label-driven, but only 1.3:1, so it is substantially about the reviewer.
%
% Related work to position against:
%   - Aroyo & Welty, "Truth is a Lie" -- disagreement as signal, not noise.
%     Closest predecessor; we add the two-kind decomposition and a policy setting.
%   - Gordon et al., jury learning -- who should decide when annotators disagree.
%   - Project Sidewalk / LabelAId -- the platform and its quality pipeline.
%   - accessibility data & municipal decision-making -- for the stakes.
%
% Venue: HCOMP or CSCW fit the crowdsourcing contribution; an urban computing
% venue works if the accessibility-policy framing leads.
% ---------------------------------------------------------------------------
